%!TEX TS-program = xelatex
\documentclass[a3paper]{adcv_color}
\usepackage[english]{babel}
\usepackage{multicol}
\usepackage{scrextend}
\setlength{\columnsep}{1cm}
\newcommand*{\tuv}{\href{https://www2.acom.ucar.edu/modeling/tropospheric-ultraviolet-and-visible-tuv-radiation-model}{\textbf{TUV }}}
\newcommand*{\aeronet}{\href{https://aeronet.gsfc.nasa.gov/}{\textbf{AERONET}}}
\newcommand*{\sedema}{\href{https://www.sedema.cdmx.gob.mx/}{\textbf{SEDEMA}~}}
\newcommand*{\omi}{\href{https://aura.gsfc.nasa.gov/omi.html}{\textbf{OMI}~}}

\title{Gamaliel-Lopez CV}
\pagestyle{empty}
\adcvname{López Padilla}{Giovanni Gamaliel}{}
\adcvemail{giovannilopez9808}{gmail}{com}
\adcvphone{(+52) 871-278-37-70}
\adcvwebsite{http://resumegglp.herokuapp.com/}{resumegglp.herokuapp.com}
%\adcvdate{}

\begin{document}
\changefontsizes{13.5pt}
Me gustan las ciencias y me gusta abordarlas desde un enfoque en el cual los problemas se puedan resolver de una manera optima mediante la implementación de la 
informática para lograr una mejor comprensión de los fenomenos. De igual manera siento una enorme curiosidad por las cosas que me rodean por lo que a veces llego a desarmar 
para ver como funciona en su interior, y asi ver que procedimientos o pasos se pueden mejorar.
\begin{multicols}{2}
\section{Experiencia en informática}

\begin{minipage}{0.8\linewidth}
  \textbf{Estimación de los tiempos de exposición solar para la obtención de la Pre vitamina D en Rosario, Argentina}\\
  Pendiente por publicar
\end{minipage}
  \begin{minipage}{0.2\linewidth}
    \begin{flushright}
    \vspace{-0.8cm}
    diciembre 2020
    \end{flushright}
  \end{minipage}\\

  Con el uso del modelo \tuv y mediciones de la columnas de ozono proporcionado por el instrumento \omi de la NASA, se creó un código en python, el cual
  consultaba los archivos HDF de \omi y creaba un archivo resultante el cual leía el modelo \tuv para así obtener los datos de irradiancia eritémica
   en la zona de Rosario, Argentina, con otro código se leían estos resultados y calculaba el tiempo de exposición solar para 
  la producción de la pre vitamina D en el cuerpo humano, tomando en cuenta los límites de exposición dependiendo el fototipo.\\

  \begin{minipage}{0.8\linewidth}
    \textbf{Análisis en la tendencia del índica UV en la Ciudad de México en el periodo 2000-2019}\\ 
    Pendiente por publicar
  \end{minipage}
    \begin{minipage}{0.2\linewidth}
      \begin{flushright}
        \vspace{-0.7cm}
      octubre\\ 2020
      \end{flushright}
    \end{minipage}\\
    
    Haciendo uso del historial de mediciones de las estaciones meteorológicas de la Secretaria del Medio Ambiente (\sedema) de la Ciudad de México,
  datos obtenidos de la red robotica de aerosoles (\aeronet) y del instrumento \omi de la NASA se realizó un estudio en la tendencia de componentes atmosféricos
  en el periodo 2000-2019, encontrando resultados acerca de las medidas que se han tomado en la reducción de la contamimación atmosférica en la región.\\

  \begin{minipage}{0.8\linewidth}
    \textbf{Localización de incendios en el Delta del río Paraná}\\
    \href{https://github.com/giovannilopez9808/Posters_templates/blob/master/Informe\%20IFIR/Focos\%20de\%20incendio\%20en\%20las\%20islas\%20del\%20Paran\%C3\%A1\%20\%20frente\%20a\%20Rosario\%20\%20(Piacentini\%2C\%20Ipi\%C3\%B1a\%2C\%20Lopez-Padilla\%2C\%20Bolmaro)\%20(Piacentini).pdf}{\textbf{Informe entregado al gobierno de Rosario, Argentina}}
  \end{minipage}
    \begin{minipage}{0.2\linewidth}
      \vspace{-0.7cm}
      \begin{flushright}
      agosto \\2020
      \end{flushright}
    \end{minipage}\\
  
  Por medio de datos satelitales proporcionados por la plataforma \href{https://firms.modaps.eosdis.nasa.gov/}{\textbf{FIRMS}} de la NASA y la NOA
  se obtiene la información sobre el número de incendios por medio de tres instrumentos satelitales: Terra, Aqua y VIIRS. Basados en estos datos satelitales se creó un código
  en Python para procesar y contabilizar el número de incendios que afectaron la región de Rosario y zonar aledañas, en el periodo del 3 de junio al 4 de agosto 2020.\\

  \begin{minipage}{0.8\linewidth}
    \textbf{Cálculo de los tiempos de exposición solar para el tratamiendo de Psoriasis en la Ciudad de México}\\
    \href{https://github.com/giovannilopez9808/Posters_templates/blob/master/2020/CNF/TES/main.pdf}{\textbf{Poster}}, \href{http://tes-v1.herokuapp.com/}{\textbf{Plataforma}}
  \end{minipage}
    \begin{minipage}{0.2\linewidth}
      \vspace{-0.9cm}
      \begin{flushright}
      agosto \\2020
      \end{flushright}
    \end{minipage}\\

  A partir de las mediciones de radiación solar eritemica proporcionadas por la Secretaria del Medio Ambiente de la 
  Ciudad de México (\sedema) se calculó el tiempo de exposición solar (TES) tomando en cuenta las caracteristicas para  diferentes fototipos, diferentes fechas y condiciones de cielo.
   Para esto se realizo un código el cual realizá la lectura de las mediciones y cálculos de los TES, en conjunto a el Centro Dermatológico de Pascua 
  se creó una plataforma llamada \textit{TES para tu salud} en donde los pacientes puedan consultar el tiempo que ellos necesitan ingresando los parámetros antes especificados.\\

  \begin{minipage}{0.8\linewidth}
    \textbf{Análisis de las irradiancias UVA y eritémica medidas por el Sistema Monitoreo atmosférico de la Ciudad de México}\\
    \href{https://github.com/giovannilopez9808/Posters_templates/blob/master/2019/AFA/Analisis\%20indice\%20UV/Analisis\%20de\%20irradiancia.pdf}{\textbf{Poster}}
  \end{minipage}
    \begin{minipage}{0.2\linewidth}
      \vspace{-0.9cm}
      \begin{flushright}
      octubre \\2019
      \end{flushright}
    \end{minipage}\\

  Con las mediciones satelitales realizadas por el instrumento \omi de la NASA, la red robotica de aerosoles (\aeronet)
  se verificó que tan aproximadas estan a la realidad los resultados del modelo \tuv, el cual fue modificado en su código fuente para que pudiera ser automatizado, ya que
  originalmente solo puede correr para un día con solo una entrada de datos.\\

  \begin{minipage}{0.8\linewidth}
    \textbf{Estudio de la irradiancia solar con dos modelos de transferencia radiativa}\\
    \href{https://github.com/giovannilopez9808/Posters_templates/tree/master/2019/CNF/Transferencia\%20radiativa}{\textbf{Poster}}
  \end{minipage}
    \begin{minipage}{0.2\linewidth}
      \vspace{-0.9cm}
      \begin{flushright}
      agosto \\2019
      \end{flushright}
    \end{minipage}\\

  Haciendo el uso del modelo \tuv y el
  modelo \href{https://www.nrel.gov/grid/solar-resource/smarts.html}{\textbf{SMARTS}}, los dos escritos originalmente en el lenguaje Fortran, se analizaron las diferencias relatividas en
  resultados para la zona metropolitana de Monterrey con los mismos parámetros. Se modificaron los códigos fuente de los dos modelos para que leyera los valores de los parámetros 
  a partir de una base de datos y así automatizar el cálculo de las irradiancias solares.
  \end{multicols}
  \begin{multicols}{2}
\section{Educación}

\textbf{Licenciatura en física}, Universidad Autónoma de Nuevo León

\section{Habilidades}

\begin{minipage}{0.25\linewidth}
  \begin{flushright}
    \textbf{Programación}\\
    \textbf{Librerías}\\
    \textbf{Hardware}\\
    \textbf{Web}\\
    \textbf{Idiomas}
  \end{flushright}
\end{minipage}
\hspace{0.3cm}
\begin{minipage}{0.75\linewidth}
  \vspace{0.25cm}
  Python, Fortran, LaTeX, Java, Julia\\
  Numpy, Matplotlib, Pandas, Scipy, Flask\\
  Arduino\\
  CSS, HTML5\\
  Español, ingles
\end{minipage}


\section{Cursos}

\begin{minipage}{0.7\linewidth}
  \textbf{Qubit x Qubit semestre 1}\\
  \textbf{Ecuaciones en diferenciales parciales desde la teoría a las aproximaciones numéricas}\\
  \textbf{Qiskit Globar Summer School, IBM}
\end{minipage}
\begin{minipage}{0.3\linewidth}
  %\vspace{-0.cm}
  \begin{flushright}
  agosto 2020\\
  agosto 2020\vspace{1cm}\\

  julio 2020
\end{flushright}
\end{minipage}
\end{multicols}
\end{document}
